%% 第五章

\chapter{实验设计与结果分析}
\chaptermark{实验设计与结果分析}
\section{实验设置}
\subsection{实验环境与数据集}
本文实验围绕面向深空场景的三维重建与新视角合成任务展开。实验输入包括多视角图像、COLMAP估计得到的相机内外参数和初始稀疏点云。训练阶段利用训练视角优化三维高斯场景表示，测试阶段在未参与训练的新视角上进行渲染，并与真实图像进行定量和定性比较。

实验数据选自M3arsSynth数据集的\textnormal{train\_video}划分。该数据集由Martian World Models工作提出，面向可控火星场景视频生成任务构建，其原始数据来源于NASA Planetary Data System(PDS)中的真实火星导航图像，并通过三维重建与渲染流程生成多视角视频序列\upcite{ref10}。本文共使用593个视频场景，其中416个场景用于整体新视角合成质量对比和模块消融实验，177个场景用于构造光照扰动数据，并训练与评估外观一致性分支。不同场景在相机轨迹、尺度范围、场景纹理丰富程度和远区域占比上存在差异，可以用于检验方法在复杂深空场景中的适应性。

对于每个视频场景，本文首先从视频序列中抽取多视角图像，并运行COLMAP完成特征提取、图像匹配、相机位姿估计和稀疏点云重建。COLMAP输出的相机参数和初始化点云随后被统一用于3DGS及各对比方法训练，以保证不同方法具有一致的数据输入和几何初始化条件。辅助深度图由Metric3Dv2-Large生成\upcite{ref30}。Metric3Dv2-Large是面向零样本度量深度和表面法向估计的单目几何基础模型，通过统一相机空间变换和深度-法向联合优化提升跨数据集泛化能力。本文不将其深度输出作为真实几何监督，而是作为相对深度排序依据，主要用于FAME远区域\textnormal{mask}构造和远区域评价。测试视角划分采用固定间隔留出策略，即每8张图像选取1张作为测试视角，其余图像作为训练视角。该设置能够在保持训练测试划分稳定的同时，评价模型对未观测视角的渲染泛化能力。实验硬件与软件环境配置如表5-1所示。

\begin{table}[H]
    \fontsize{10.5pt}{13.125pt}\selectfont
    \caption{实验硬件与软件环境配置}
    \centering\label{tab:5-1}
    \begin{adjustbox}{max width=\textwidth}
    \begin{tabular}{cc}
        \toprule
        环境项目 & 配置说明 \\
        \midrule
        操作系统 & Ubuntu 24.04.2 LTS \\
        处理器 & Intel(R) Xeon(R) Silver 4210R CPU @ 2.40GHz \\
        显卡 & NVIDIA GeForce RTX 3090，24576 MiB \\
        CUDA 版本 & 11.6 \\
        Python 版本 & 3.7.12 \\
        PyTorch 版本 & 1.12.1 \\
        COLMAP 版本 & 3.14.0.dev0 \\
        \bottomrule
    \end{tabular}
    \end{adjustbox}
\end{table}
\subsection{对比方法与实验设置}
为验证本文方法的有效性，实验选取基础3D Gaussian Splatting作为主要基线，并与D2GS、3D-HGS、SCGaussian、ERank-GS和MVGS等相关方法进行对比。针对光照扰动条件下的外观一致性实验，进一步选取Luminance-GS和WildGaussians作为外观变化建模方向的对比方法。本文方法记为Ours。

所有方法尽量采用相同的数据输入、训练测试划分和评价脚本。本文方法在3DGS框架上加入FAME远区域增强、有效秩正则、体积正则和外观建模相关约束。训练过程中远区域划分比例采用第四章确定的0.30，三维高斯模型在统一迭代次数下保存并用于测试视角渲染。对于光照扰动实验，仅改变训练图像的光照与颜色条件，不改变场景几何、相机位姿和测试视角划分，以便单独分析外观变化对新视角合成结果的影响。

\subsection{实验协议与公平性说明}
为保证不同方法之间的实验结果具有可比性，本文在实验过程中尽量保持数据输入、训练测试划分、相机参数、渲染视角和评价指标的一致性。对于每个场景，所有方法均使用相同的多视角图像、COLMAP相机参数和测试视角划分。训练视角用于三维高斯场景表示的优化，测试视角不参与训练过程，仅用于新视角渲染和指标计算。通过固定训练测试划分，可以避免由于测试视角不同导致的结果偏差，使不同方法之间的对比更加公平。

此外，为减少批量实验过程中人为筛选带来的影响，本文在结果统计时采用统一的场景遍历、渲染和指标计算流程。每个场景完成训练后，均按照固定迭代次数加载模型，并在相同测试视角下生成渲染结果；随后由统一评价脚本自动读取render图像、gt图像以及对应的评价区域，计算各项指标并汇总平均结果。对于不同方法产生的输出目录，本文保持一致的结果组织方式，使指标统计过程不依赖人工挑选图像或人工判断局部效果。通过这种批量化、自动化的评价方式，可以降低主观选择对实验结论的影响，使最终结果更能反映不同方法在整体数据集上的平均表现。

在对比方法设置方面，本文尽量保持各方法原始实现中的核心参数和默认训练策略，同时保证输入数据和评价方式一致。由于不同方法的代码结构、训练策略和适用场景存在差异，部分方法可能具有特定的初始化方式、正则项或训练超参数。本文不对这些方法进行过度改写，而是在统一数据和统一评价协议下运行，以尽量保留各方法本身的设计特点。对于无法稳定适配本文数据或在部分场景中运行失败的方法，本文在统计时仅使用能够正常完成训练、渲染和评价流程的有效结果，并在结果分析中结合方法特点进行解释。

\subsection{默认超参数设置}
本文默认超参数设置如下：基础重建损失中的SSIM权重$\lambda_{\mathrm{ssim}}=0.2$；FAME模块的远区域比例$r_{\mathrm{far}}=0.30$，损失权重$\lambda_{\mathrm{fame}}=1.0$；RVS模块中的有效秩正则权重$\lambda_{\mathrm{erank}}=0.05$，体积正则权重$\lambda_{\mathrm{vol}}=0.01$；启用外观一致性分支时，采样坐标正则权重$\lambda_{\mathrm{sc}}=0.01$。所有方法使用相同的训练测试划分、相机参数、初始点云和评价流程；除本文新增模块对应参数外，对比方法尽量保持其公开实现或原始配置中的默认训练设置。为保证不同实验之间具有可比性，整体新视角合成对比、远区域分析和消融实验均采用上述默认设置；光照扰动实验仅改变训练图像外观条件，不改变场景几何、相机位姿和测试视角划分。

\section{新视角合成质量对比}
新视角合成质量对比是本文实验的主要结果，用于评价不同方法在测试视角下的整体图像重建能力。表5-2给出了不同方法在深空场景上的PSNR、SSIM和LPIPS指标结果。

\begin{table}[H]
    \fontsize{10.5pt}{13.125pt}\selectfont
    \caption{不同方法在深空场景上的新视角合成质量对比}
    \centering\label{tab:5-2}
    \begin{adjustbox}{max width=\textwidth}
    \begin{tabular}{cccc}
        \toprule
        Method & SSIM$\uparrow$ & PSNR$\uparrow$ & LPIPS$\downarrow$ \\
        \midrule
        3D Gaussian Splatting (ACM TOG’23) & 0.931 & 36.17 & 0.131 \\
        D$^2$GS (ICLR’26) & 0.951 & 35.31 & 0.086 \\
        3D-HGS (CVPR’25) & 0.927 & 36.16 & 0.136 \\
        SCGaussian (NeurIPS’24) & 0.535 & 21.91 & 0.548 \\
        ERank-GS (NeurIPS’24) & 0.931 & 35.85 & 0.129 \\
        MVGS (arXiv’24) & 0.956 & 38.27 & 0.089 \\
        \textbf{Ours} & \textbf{0.961} & \textbf{38.49} & \textbf{0.078} \\
        \bottomrule
    \end{tabular}
    \end{adjustbox}
\end{table}
从不同方法的表现来看，基础3D Gaussian Splatting在本文数据集上已经能够取得较高的整体PSNR和SSIM，说明显式三维高斯表示适合用于深空场景的新视角合成任务。但是，3DGS主要依赖全图图像重建损失进行优化，在远区域纹理弱、场景结构尺度变化大以及视角约束不足的情况下，容易优先拟合近处纹理明显区域，而对远处区域层次、小尺度结构和弱纹理区域的关注不足。因此，虽然3DGS整体指标较高，但在远区域细节和局部结构稳定性方面仍存在进一步提升空间。

D$^2$GS从深度和密度角度增强稀疏视角条件下的重建稳定性，其SSIM和LPIPS相较3DGS有明显改善，说明深度先验和密度约束能够提升结构保持和感知质量。然而，D$^2$GS的PSNR低于3DGS，表明深度或密度引导虽然有助于缓解几何退化，但并不一定能够同步提升像素级重建精度。对于本文关注的深空场景而言，辅助深度信息本身可能受到弱纹理、阴影和远距离成像的影响，因此如何利用深度先验而不过度依赖其局部误差，是方法设计中需要平衡的问题。

3D-HGS通过半高斯结构增强表面表达能力，在一些表面重建任务中具有较强的几何动机。但在本文实验结果中，3D-HGS与3DGS相比提升并不明显，PSNR基本接近，而SSIM和LPIPS略低。这说明改变高斯基元形态并不必然带来整体新视角合成质量提升。深空场景通常存在大范围弱纹理区域和远近尺度差异，仅依靠更强的表面基元假设，可能难以充分解决远区域监督不足和跨视角外观不一致问题。

SCGaussian在本文数据集上的整体表现相对较弱。该方法强调利用结构一致性和匹配先验改善少样本新视角合成，但本文数据中存在大量低纹理、重复纹理和远距离场景区域，这些因素容易削弱匹配先验的可靠性。当匹配关系本身不稳定时，基于匹配的结构约束可能无法有效发挥作用，甚至会影响后续三维高斯优化。因此，该结果并不意味着SCGaussian方法本身无效，而是说明其依赖的结构匹配条件与本文深空数据的成像特点并不完全一致。

ERank-GS在本文数据集上的结果与3DGS较为接近，PSNR略低于3DGS，SSIM基本保持一致，LPIPS也只有小幅变化。这说明该方法在本文深空场景中的整体收益相对有限。造成这一现象的可能原因在于，本文数据中存在大量远距离弱纹理区域、重复场景纹理和明显尺度差异，而ERank-GS的改进重点并不是专门针对远区域监督不足或跨视角外观变化设计的。因此，在本文任务设定下，它虽然能够保持较稳定的基础渲染能力，但对远处细节恢复和整体感知质量的提升不如本文方法明显。相比之下，本文方法进一步结合远区域增强和外观一致性处理，更贴合深空场景中的主要退化来源。

MVGS在本文实验中表现较强，PSNR、SSIM和LPIPS均明显优于3DGS，说明多视角约束和跨视角信息利用对于提升深空场景重建质量具有重要作用。与MVGS相比，本文方法在三项指标上仍取得进一步提升，说明远区域增强、高斯形态稳定约束和外观建模能够在多视图约束之外进一步改善新视角合成效果。尤其对于远处区域、弱纹理区域和局部边缘结构，本文方法能够通过FAME增强监督信号，并通过RVS抑制异常高斯形态，从而获得更加稳定的渲染结果。

为进一步分析本文方法在不同难度场景下的稳定性，本文根据基础3DGS在各场景上的PSNR指标对全部测试场景进行排序，并按照排序结果划分为困难组、常规组和简单组。其中，PSNR较低的前15\%场景划分为困难组，PSNR较高的后15\%场景划分为简单组，中间70\%场景划分为常规组。该划分方式以基础3DGS的重建结果作为难度依据，能够反映不同场景在原始高斯表示下的重建难易程度。通常情况下，困难组场景更可能包含远区域欠拟合、弱纹理、视角约束不足等问题；简单组场景则表示基础3DGS已经能够取得较好重建效果的场景。

表5-3给出了不同难度场景下本文方法相对3DGS的性能提升结果。为避免仅通过全体场景平均值掩盖不同难度场景中的性能差异，本文分别统计了各难度组中3DGS与Ours的PSNR、SSIM和LPIPS平均值，并计算二者之间的变化量。

\begin{table}[H]
    \fontsize{10.5pt}{13.125pt}\selectfont
    \caption{不同难度场景下本文方法相对3DGS的性能提升对比}
    \centering\label{tab:5-3}
    \begin{adjustbox}{max width=\textwidth}
    \begin{tabular}{ccccccccccc}
        \toprule
        难度分组 & 场景数量 & 3DGS PSNR & Ours PSNR & $\Delta$PSNR & 3DGS SSIM & Ours SSIM & $\Delta$SSIM & 3DGS LPIPS & Ours LPIPS & $\Delta$LPIPS \\
        \midrule
        困难组 & 62 & 27.62 & 33.81 & 5.70 & 0.786 & 0.896 & 0.088 & 0.355 & 0.200 & 0.134 \\
        常规组 & 292 & 36.99 & 38.87 & 1.60 & 0.951 & 0.971 & 0.015 & 0.101 & 0.061 & 0.035 \\
        简单组 & 62 & 40.90 & 41.33 & 0.36 & 0.980 & 0.982 & 0.002 & 0.046 & 0.036 & 0.007 \\
        全部场景 & 416 & 36.17 & 38.49 & 2.31 & 0.931 & 0.961 & 0.030 & 0.131 & 0.078 & 0.053 \\
        \bottomrule
    \end{tabular}
    \end{adjustbox}
\end{table}
由表5-3可以看出，本文方法在不同难度场景下均优于基础3DGS，且在困难组中的提升最为明显。在困难组中，Ours的PSNR由3DGS的27.62dB提升至33.81dB，提升5.70dB；SSIM由0.786提升至0.896，提升0.088；LPIPS由0.355降低至0.200，降低0.134。该结果说明，在基础3DGS难以充分重建的场景中，本文方法能够更有效地改善图像重建质量和感知质量。

在常规组中，本文方法同样保持稳定提升，PSNR提升1.60dB，SSIM提升0.015，LPIPS降低0.035。这说明本文方法并非仅对少数困难样本有效，而是在占比较高的常规场景中也能够持续改善新视角合成结果。对于简单组，基础3DGS已经取得较高指标，PSNR达到40.90dB，SSIM达到0.980，LPIPS仅为0.046，因此进一步提升空间相对有限。在这种情况下，本文方法仍将PSNR提升至41.33dB，并将LPIPS降低至0.036，说明本文方法在简单场景中没有破坏原有较好的重建结果，而是仍然能够带来小幅稳定增益。

综合不同难度组结果可以发现，本文方法的提升幅度与场景难度具有一定相关性：场景越困难，本文方法相对3DGS的提升越明显。这一现象与本文的设计目标一致。困难场景中往往存在远区域纹理弱、局部结构尺度变化大、视角约束不足或外观变化明显等问题，基础3DGS仅依赖全图重建损失时容易出现远处细节弱化、局部结构模糊或空间表达不稳定现象。本文通过远区域增强约束提高远处区域在训练过程中的影响，同时通过高斯形态稳定约束缓解尺度退化和异常拉伸问题，从而在基础方法表现较弱的场景中取得更明显的改进。该结果进一步说明，本文方法不仅能够提升整体平均指标，而且能够重点改善深空场景中更具挑战性的重建情况。

\begin{figure}[H]
    \centering
    \includegraphics[width=0.78\textwidth]{深空场景新视角合成可视化对比.jpg}
    \caption{深空场景新视角合成可视化对比}\label{fig:5-1}
\end{figure}
从图5-1可视化结果也可以看出，3DGS在远处区域和弱纹理区域容易出现局部模糊、纹理连续性不足和小尺度结构弱化等现象；相比之下，本文方法在局部颗粒、目标边缘和远处背景层次上保持更加清晰，渲染结果与真实图像更接近。该可视化结果与表5-3中PSNR、SSIM和LPIPS的提升趋势一致，说明本文方法不仅改善了整体图像指标，也提升了深空场景中容易退化区域的视觉质量。

\section{远区域增强效果分析}
为验证FAME模块对远区域细节恢复的作用，本文选取典型深空场景进行可视化对比，如图5-2所示。其中，第一列为原始3DGS方法，第二列为在3DGS基础上加入FAME模块后的结果，第三列为真实图像。红色矩形框标出了远区域中需要重点观察的局部区域，下方对应展示了红框区域的局部放大结果。

\begin{figure}[H]
    \centering
    \includegraphics[width=0.78\textwidth]{与+FAME模块对比.jpg}
    \caption{FAME模块对远区域细节恢复效果的可视化对比}\label{fig:5-2}
\end{figure}
\vspace{-\baselineskip}
\begin{figure}[H]
    \centering
    \includegraphics[width=0.56\textwidth]{局部放大图.png}
    \caption{局部放大结果对比}\label{fig:5-3}
\end{figure}
\vspace{-\baselineskip}
\begin{table}[H]
    \fontsize{10.5pt}{13.125pt}\selectfont
    \caption{不同方法在远区域上的新视角合成质量对比}
    \centering\label{tab:5-4}
    \begin{adjustbox}{max width=\textwidth}
    \begin{tabular}{cccc}
        \toprule
        Method & SSIM$\uparrow$ & PSNR$\uparrow$ & LPIPS$\downarrow$ \\
        \midrule
        mvgs & 0.943 & 37.13 & 0.134 \\
        3dhgs & 0.909 & 34.91 & 0.182 \\
        scgaussian & 0.655 & 22.95 & 0.442 \\
        3DGS & 0.906 & 35.14 & 0.180 \\
        3DGS + FAME & 0.940 & 37.73 & 0.118 \\
        \textbf{Ours} & \textbf{0.953} & \textbf{38.41} & \textbf{0.094} \\
        \bottomrule
    \end{tabular}
    \end{adjustbox}
\end{table}
从远区域评价结果可以看出，本文方法在远区域PSNR、SSIM和LPIPS上均取得较好表现，说明FAME远区域增强模块能够有效改善深空场景中远处结构的渲染质量。表5-4进一步给出了不同方法在远区域上的定量对比结果。与原始3DGS相比，加入FAME后SSIM由0.906提升至0.940，PSNR由35.14dB提升至37.73dB，LPIPS由0.180降低至0.118，表明FAME模块能够在结构相似性、像素重建误差和感知质量三个方面同时带来提升。

造成这一现象的主要原因在于，3DGS训练时通常对整幅图像计算重建损失，而远区域在图像中往往只占据较小比例。相比前景岩石、近处纹理或高对比边缘，远处区域的像素数量更少、纹理变化更弱，因而在全图平均损失中的贡献容易被削弱。模型为了快速降低整体重建误差，通常会优先拟合近处区域和纹理明显区域，而远区域则可能出现细节模糊、层次不足或边缘弱化等问题。

针对上述问题，FAME模块通过辅助深度结果构造远区域\textnormal{mask}，并在基础全图重建损失之外，额外计算远区域内的重建误差。这样，远处区域像素不再只是全图损失中的少量组成部分，而是作为一个独立的优化目标参与反向传播。该机制使模型在训练过程中更加关注远处区域纹理、背景层次和小尺度结构，从而缓解远区域被近景区域主导的问题。

从视觉表现上看，FAME对远区域的改善主要体现在三个方面。第一，远处区域纹理更加清晰，原本容易被平滑化的弱纹理区域能够保留更多亮度变化和颗粒结构。第二，远区域中岩石边缘、结构起伏和背景轮廓更加稳定，减少了模糊边界和局部断裂现象。第三，远近区域之间的层次关系更加自然，模型不容易将远处区域简单拟合为平滑背景，从而提升整体场景的空间感。

\section{RVS模块对高斯空间表示稳定性的分析}
3D Gaussian Splatting中每个高斯基元具有三个尺度方向，其尺度分布直接影响高斯在三维空间中的形态。当某一尺度方向远大于其他方向时，高斯容易退化为针状结构；当整体尺度异常增大或缩小时，则可能造成局部空间覆盖不稳定。这类异常高斯基元虽然不一定直接对应真实几何误差，但会影响局部表面表达、边缘清晰度和新视角渲染稳定性。因此，为进一步分析RVS模块对高斯基元形态的约束作用，本文从尺度比例分布和异常高斯数量两个角度进行可视化统计。

设某个高斯基元三个尺度方向中的最大值和最小值分别为$s_{\max}$和$s_{\min}$，本文采用尺度比例$s_{\max}/s_{\min}$衡量高斯形态的各向异性程度。该比例越大，说明该高斯在某一方向上的拉伸越明显；当比例超过设定阈值时，可认为该高斯存在异常尺度退化风险。本文进一步统计尺度比例分布及其累积分布函数，用于观察不同方法下高斯基元是否存在明显长尾现象。

图5-4展示了3DGS与加入RVS后的尺度比例分布对比。其中，红色虚线表示异常尺度比例阈值。可以看到，3DGS的尺度比例分布存在明显长尾，部分高斯的$s_{\max}/s_{\min}$达到较大数值，说明训练过程中出现了较多异常拉伸的高斯基元。这类高斯通常更容易出现在远区域、弱纹理区域或视角约束不足的位置，可能导致局部漂浮、边缘模糊或表面表达不连续等现象。相比之下，加入RVS后，高斯尺度比例明显向较小范围集中，超过异常阈值的高斯数量显著减少，说明有效秩正则和体积正则能够对尺度分布和空间覆盖范围形成约束，从而抑制极端尺度退化。

\begin{figure}[H]
    \centering
    \begin{subfigure}[b]{0.49\linewidth}
        \includegraphics[width=\linewidth]{30000_baseline.png}
        \caption{3DGS}
    \end{subfigure}
    \begin{subfigure}[b]{0.49\linewidth}
        \includegraphics[width=\linewidth]{30000_RVS.png}
        \caption{3DGS+RVS}
    \end{subfigure}
    \caption{RVS对高斯尺度比例分布的影响}\label{fig:5-4}
\end{figure}
进一步地，本文统计了不同测试视角下可见异常高斯的数量，如图5-5所示。3DGS在各个视角中可见的异常高斯数量均较高，而加入RVS后异常高斯数量大幅下降。以图中统计结果为例，基础设置下异常高斯总量约为90.3k，而加入RVS后降低至约810；在各测试视角中，3DGS的可见异常高斯数量约为78.0k至84.2k，而加入RVS后降低至约719至767。这说明RVS并非只在少数视角中产生作用，而是在不同观察视角下均能稳定抑制异常尺度高斯，使显式高斯表示更加规整。

\begin{figure}[H]
    \centering
    \includegraphics[width=0.92\textwidth]{异常高斯数对比.png}
    \caption{不同视角下异常高斯数量对比}\label{fig:5-5}
\end{figure}
从上述高斯尺度比例分布和异常高斯数量统计可以看出，RVS并不是直接提高某一张测试图像的像素拟合能力，而是通过约束三维高斯基元的空间形态，使场景表示更加稳定。对于3DGS而言，最终渲染图像是大量三维高斯投影到二维图像平面后进行颜色累积的结果。如果部分高斯在三维空间中过度拉伸或体积异常膨胀，其投影到图像平面后可能覆盖不合理的区域，进而造成局部颜色扩散、边界模糊或遮挡关系不稳定等问题。尤其在远区域和弱纹理区域，图像监督信号本身较弱，异常高斯更容易在优化过程中被保留下来，从而影响新视角下的结构一致性。

因此，RVS对新视角合成质量的提升具有间接性。它并不是像FAME那样直接改变远区域像素的损失权重，而是从三维表示本身出发，减少异常尺度高斯对局部空间结构的干扰。当高斯基元的尺度方向分布更加合理、整体体积变化更加稳定时，不同视角下的颜色累积过程会更加一致，局部表面支撑也会更加可靠。这也是本文方法能够在SSIM和LPIPS等结构与感知指标上取得提升的重要原因。

\section{光照扰动条件下的外观一致性分析}
光照扰动实验用于分析模型在复杂外观条件下的新视角合成鲁棒性。实验在保持场景几何、相机位姿和训练测试划分不变的前提下，对训练图像构造低光、过曝、白平衡偏移和局部亮度变化等扰动，从而模拟深空场景中可能出现的非均匀光照和成像差异。需要说明的是，光照扰动实验与常规新视角合成实验的关注重点不同。常规实验主要评价模型在原始数据条件下的整体重建质量，而光照扰动实验更关注模型在训练图像存在亮度、曝光和颜色差异时的外观适应能力。若将二者混合分析，容易将几何重建误差和外观变化误差混在一起，导致实验结论不够清晰。因此，本文在保持场景几何结构、相机位姿和测试视角划分不变的前提下，单独构造光照扰动数据，用于观察不同方法在非一致外观条件下的鲁棒性。

在该实验设置下，如果模型缺少有效的外观建模能力，就可能将不同视角之间的亮度变化错误解释为场景颜色或局部结构变化，进而导致渲染结果出现颜色偏移、局部发灰、纹理不稳定等现象。对于深空场景而言，太阳高度角、阴影范围、曝光状态和成像响应差异都可能造成类似问题。因此，光照扰动实验不仅用于验证本文外观建模分支的有效性，也用于说明深空场景中新视角合成任务并不只是几何重建问题，还需要同时考虑跨视角外观一致性。表5-5给出了不同方法在光照扰动条件下的新视角合成质量对比。

\begin{table}[H]
    \fontsize{10.5pt}{13.125pt}\selectfont
    \caption{光照扰动条件下的新视角合成质量对比}
    \centering\label{tab:5-5}
    \begin{adjustbox}{max width=\textwidth}
    \begin{tabular}{cccc}
        \toprule
        Method & SSIM$\uparrow$ & PSNR$\uparrow$ & LPIPS$\downarrow$ \\
        \midrule
        3D Gaussian Splatting (ACM TOG’23) & 0.936 & 29.76 & 0.100 \\
        Luminance-GS (CVPR’25) & 0.873 & 20.32 & 0.198 \\
        Wildgaussians (NeurIPS’24) & 0.768 & 23.38 & 0.282 \\
        \textbf{Ours} & \textbf{0.942} & \textbf{33.48} & \textbf{0.100} \\
        \bottomrule
    \end{tabular}
    \end{adjustbox}
\end{table}
由表5-5可知，在光照扰动条件下，本文方法在SSIM和PSNR上均取得最优结果，LPIPS与3DGS持平并优于其他外观建模方法。该结果说明本文外观建模设计能够在一定程度上缓解光照变化带来的颜色偏差和曝光不一致问题，使新视角渲染结果在结构保持和像素重建方面更加稳定。

\begin{figure}[H]
    \centering
    \includegraphics[width=0.82\textwidth]{外观一致性可视化结果.jpg}
    \caption{光照扰动条件下的新视角合成可视化对比}\label{fig:5-6}
\end{figure}
从图5-6可以看出，在复杂光照扰动条件下，开启外观一致性模块后的结果整体上更接近GT。第一行主要包含弱纹理区域和大范围明暗渐变区域，未开启外观一致性模块时，图像右侧暗部区域的色调偏移更加明显，场景纹理对比度也与GT存在差异；开启该模块后，整体亮度分布和暗部过渡更加稳定。第二行包含起伏结构、碎石和局部阴影区域，未开启模块时，碎石边缘和阴影附近容易出现颜色不均衡现象，而开启模块后，岩石区域的亮度层次和场景颜色更接近真实图像。第三行包含较多小尺度石块和局部颗粒结构，未开启模块时局部色调略有漂移，前景区域的亮暗关系不够自然；开启模块后，场景颜色、阴影边界以及小石块周围的外观变化更加连续。总体而言，外观一致性模块能够在低光、过曝、白平衡偏移和局部亮度变化等扰动下，提升模型对颜色分布和局部外观变化的表达能力，从而增强新视角合成结果的视觉一致性。

\section{消融实验分析}
消融实验用于分析本文不同模块对新视角合成质量的贡献。实验以3DGS为基线，分别加入有效秩正则、体积正则、RVS、FAME以及FAME与RVS的组合，比较不同配置下的PSNR、SSIM和LPIPS指标变化。其中，RVS表示有效秩正则与体积正则的联合设置，主要用于约束高斯基元的尺度分布和整体体积变化；FAME主要作用于远区域图像重建监督。二者分别从高斯形态稳定性和远区域监督增强两个角度对3DGS进行补充。为展示各组件是否启用对指标的影响，表5-6以勾选形式列出了有效秩正则、体积正则、RVS和FAME的组合关系及对应的新视角合成质量。

\begin{table}[H]
    \fontsize{10.5pt}{13.125pt}\selectfont
    \caption{本文方法各组件对新视角合成质量的影响}
    \centering\label{tab:5-6}
    \begin{adjustbox}{max width=\textwidth}
    \begin{tabular}{ccccccc}
        \toprule
        Effective Rank & Volume Reg. & RVS & FAME & PSNR($\uparrow$) & SSIM($\uparrow$) & LPIPS($\downarrow$) \\
        \midrule
         &  &  &  & 36.17 & 0.931 & 0.131 \\
        \checkmark &  &  &  & 36.91 & 0.945 & 0.111 \\
         & \checkmark &  &  & 36.30 & 0.931 & 0.130 \\
         &  & \checkmark &  & 37.64 & 0.954 & 0.093 \\
         &  &  & \checkmark & 37.69 & 0.945 & 0.104 \\
         &  & \textbf{\checkmark} & \textbf{\checkmark} & \textbf{38.49} & \textbf{0.961} & \textbf{0.078} \\
        \bottomrule
    \end{tabular}
    \end{adjustbox}
\end{table}
注：RVS表示同时启用有效秩正则与体积正则后的联合设置。表中单独列出Effective Rank和Volume Reg.是为了观察二者单独作用，RVS行用于表示二者联合后的整体效果。

从表5-6可以看出，不同模块对新视角合成质量的影响侧重点并不完全相同。与基础3DGS相比，单独加入有效秩正则后，SSIM由0.931提升至0.945，LPIPS由0.131降低至0.111，说明该约束能够明显改善渲染结果的结构保持能力和感知质量。其原因在于，有效秩正则主要关注高斯三个尺度方向之间的相对分布，能够减少单一方向过度拉伸导致的局部结构不稳定问题，使高斯基元在空间中形成更加合理的支撑形态。因此，该模块虽然不是直接作用于图像颜色本身，但能够通过改善三维表示的局部稳定性，间接提升新视角渲染质量。

单独加入体积正则时，PSNR由36.17dB小幅提升至36.30dB，LPIPS由0.131降低至0.130，整体提升相对有限。这说明仅约束高斯整体体积不足以显著改善渲染质量。体积正则主要限制高斯基元的整体空间覆盖范围，能够抑制局部高斯过度膨胀或过度收缩，但它并不能充分约束三个尺度方向之间的比例关系。也就是说，即使某个高斯的整体体积处于合理范围内，它仍然可能存在单一方向过长或局部形态不均衡的问题。因此，体积正则单独使用时对渲染质量的提升较弱，更多起到辅助约束作用。

当有效秩正则与体积正则联合形成RVS后，PSNR、SSIM和LPIPS分别达到37.64dB、0.954和0.093，明显优于二者单独使用的结果。这表明尺度方向约束与体积约束具有互补作用：有效秩正则主要抑制方向性退化，使高斯三个尺度方向之间的分布更加均衡；体积正则主要限制整体尺度异常，避免高斯过度膨胀或过度收缩。二者结合后，既能够改善高斯基元的形态比例，又能够控制其空间覆盖范围，从而使显式高斯表示更加稳定，并进一步提升新视角合成质量。

结合前文高斯形态稳定性分析可以看出，RVS的作用是通过有效秩正则和体积正则减少异常尺度高斯，使高斯基元的尺度分布更加规整。高斯基元形态更加稳定后，局部区域的颜色累积、遮挡关系和表面支撑过程更加可靠，因此能够间接提升渲染结果的结构一致性和感知质量。这也解释了为什么RVS在SSIM和LPIPS上带来较明显改善：SSIM更关注结构相似性，LPIPS更关注感知差异，而高斯基元空间分布的稳定性正会影响局部结构和视觉细节的表达效果。

相比之下，FAME主要通过增强远区域监督改善远处区域纹理和目标细节。该模块并不直接约束高斯基元的尺度形态，而是通过远区域\textnormal{mask}提高远处像素在训练损失中的影响力，使模型更加关注远处区域、背景层次和小尺度结构。因此，FAME更直接地改善了远区域细节恢复问题，并能够在整体PSNR上带来明显增益。也就是说，FAME主要解决的是“远区域监督不足”的问题，而RVS主要解决的是“显式高斯表示不稳定”的问题，二者的作用层面并不相同。

当FAME与RVS同时使用时，PSNR、SSIM和LPIPS均达到最优，说明远区域监督增强与高斯形态稳定约束具有互补关系。FAME从图像监督层面增强远区域重建信号，使模型能够更充分地学习远处区域纹理和目标结构；RVS从三维表示层面约束高斯基元尺度分布，减少异常拉伸、局部膨胀和不稳定表示。二者结合后，模型既能改善远区域细节，又能获得更加稳定的显式高斯表示，因此在像素重建误差、结构相似性和感知质量三个方面均取得更优结果。整体消融结果说明，本文方法的性能提升并不是由单一模块带来的，而是来自远区域监督增强和高斯形态约束之间的协同作用。

\section{本章小结}
本章围绕深空场景的新视角合成任务，对本文方法进行了实验验证。首先介绍了实验环境、数据集、对比方法和评价指标，并采用PSNR、SSIM和LPIPS作为主要图像质量评价标准。随后，通过整体新视角合成质量对比验证了本文方法相较基础3DGS和其他对比方法的优势；通过远区域实验分析了FAME对远处区域纹理、目标边缘和背景层次的增强作用；通过RVS相关实验说明高斯空间形态稳定约束能够改善显式高斯表示的稳定性；通过光照扰动实验验证了外观建模策略在复杂成像条件下的有效性；最后通过消融实验进一步分析了不同模块对最终性能的贡献。实验结果表明，本文方法能够有效提升深空场景的新视角合成质量，并在远区域渲染和复杂外观条件下表现出更好的适应能力。
